\chapter{Physical Meaning of the Effective Action and Effective Potential}
\label{ch:physical-meaning}

We have defined $\Gamma[\varphi_c]$ and $V_{\text{eff}}(\bar\varphi)$
formally in the previous chapter. Here we give them concrete physical
meaning.

\section{The Effective Potential as a Generator of 1PI Graphs}

For a constant classical field $\bar{\varphi}$, inserting the Fourier
representation of $\Gamma^{(n)}$ into the expansion
\eqref{eq:gamma-1PI-expansion} gives
\begin{align}
	\Gamma(\bar{\varphi}) & = \sum_{n=2}^{\infty} \frac{1}{n!}\,
	\bar{\varphi}^n \int \prod_{i=1}^{n} d^4x_i\;
	\Gamma^{(n)}(x_1, \ldots, x_n)                                               \\
	                      & = \sum_{n=2}^{\infty} \frac{1}{n!}\, \bar{\varphi}^n
	\int \prod_{i=1}^{n} \left( d^4x_i\, \frac{d^4k_i}{(2\pi)^4}\,
	e^{ik_i x_i} \right) \tilde{\Gamma}^{(n)}(k_1, \ldots, k_n) \cdot
	(2\pi)^4\, \delta^{(4)}\!\left(\sum_{i=1}^{n} k_i\right),
\end{align}
where $\int d^4x\, e^{ikx} = (2\pi)^4\delta^{(4)}(k)$ was used. Setting
$k=0$ gives $\int d^4x = (2\pi)^4\delta^{(4)}(0)$, so that
\begin{equation}
	\Gamma(\bar{\varphi}) = \sum_{n=2}^{\infty} \frac{1}{n!}\,
	\bar{\varphi}^n \int d^4x\; \tilde{\Gamma}^{(n)}(0, \ldots, 0).
\end{equation}
Comparing this with $\Gamma(\bar{\varphi}) =
	\int d^4x\,[-V_{\text{eff}}(\bar{\varphi})]$ gives
\begin{equation}
	V_{\text{eff}}(\bar{\varphi}) = -\sum_{n=2}^{\infty}
	\frac{1}{n!}\,\bar{\varphi}^n\,\tilde{\Gamma}^{(n)}(0, \ldots, 0).
\end{equation}
Thus $V_{\text{eff}}(\bar\varphi)$ is the generating functional for 1PI
graphs at vanishing external momenta, computed by summing all such graphs
and inserting a factor $\bar\varphi$ for each external line.

To lowest order in $\hbar$, only tree-level graphs contribute. For a
Lagrangian of the form
\begin{equation}
	\mathcal{L}(\varphi, \partial_\mu \varphi) = \frac{1}{2}\,\partial_\mu
	\varphi\, \partial^\mu \varphi - U(\varphi),
\end{equation}
each term $g_i\varphi^i/i!$ in $U(\varphi)$ contributes exactly one
tree-level 1PI graph, since all tree graphs with more than one vertex are
not 1PI. Each term gives a contribution $-g_i\bar\varphi^i$, where $-g_i$
is the vertex factor and $\bar\varphi^i$ comes from the $i$ external
lines. The minus signs cancel and one recovers
\begin{equation}
	V_{\text{eff}}(\bar{\varphi}) = U(\bar{\varphi}) + \mathcal{O}(\hbar),
\end{equation}
confirming that the classical potential is the leading-order approximation
to the effective potential.

\section{The Effective Potential as a Minimum Energy Density}

A deeper interpretation follows from the variational principle.

\begin{theorem}
	$V_{\mathrm{eff}}(\bar{\varphi})$ is the minimum of the energy density
	expectation value over all normalized states $|a\rangle$ satisfying
	$\langle a|\varphi|a\rangle = \bar{\varphi}$. That is,
	\begin{equation}
		V_{\text{eff}}(\bar{\varphi}) = \langle a|\mathcal{H}|a\rangle,
	\end{equation}
	for $|a\rangle$ which obeys $\delta\langle a|\mathcal{H}|a\rangle = 0$
	subject to the constraints $\langle a|a\rangle = 1$,
	$\langle a|\varphi|a\rangle = \bar{\varphi}$.
\end{theorem}

\begin{proof}
	We first solve the analogous problem in quantum mechanics, then translate
	to field theory. In QM the problem is to find a state $|a\rangle$ such
	that $\langle a|H|a\rangle$ is stationary subject to $\langle a|a\rangle
		= 1$ and $\langle a|A|a\rangle = A_c$ for some self-adjoint operator $A$.
	Introducing Lagrange multipliers $E$ and $J$ for each constraint, the
	problem becomes
	\begin{align}
		\delta\langle a|(H - E - JA)|a\rangle & = 0 \quad \text{(unconstrained)} \\
		\langle a|a\rangle                    & = 1                              \\
		\langle a|A|a\rangle                  & = A_c.
	\end{align}
	Since $H$ is self-adjoint, this yields
	\begin{equation}
		(H - E - JA)|a\rangle = 0,
	\end{equation}
	with solution $|a\rangle = |a(E,J)\rangle$. Using $\langle a|a\rangle=1$,
	$E$ is solved as a function of $J$, i.e.\ $E = E(J)$, and
	$\langle a|A|a\rangle = A_c$ gives $J = J(A_c)$. Thus
	\begin{equation}
		(H - JA)|a(J)\rangle = E(J)|a(J)\rangle,
	\end{equation}
	where $|a(J)\rangle$ is a normalized eigenstate of $(H-JA)$. Taking the
	expectation value with $\langle a(J)|$ and differentiating with respect
	to $J$,
	\begin{align}
		\left(\frac{\delta}{\delta J}\langle a(J)|\right)
		\underbrace{(H-E(J)-JA)|a(J)\rangle}_{0}
		 & + \langle a(J)|\frac{\delta}{\delta J}
		(H - E(J) - JA)|a(J)\rangle \notag        \\
		 & + \langle a(J)|(H-E(J)-JA)
		\underbrace{\frac{\delta}{\delta J}|a(J)\rangle}_{0} = 0,
	\end{align}
	which gives
	\begin{equation}
		-\frac{\delta E}{\delta J} -
		\underbrace{\langle a(J)|A|a(J)\rangle}_{A_c} = 0
		\qquad \Rightarrow \qquad
		\frac{\delta E}{\delta J} = -A_c.
	\end{equation}
	Substituting back into the expectation value of the eigenvalue equation,
	\begin{equation}
		\langle a(J)|H|a(J)\rangle = E(J) - J\frac{\delta E}{\delta J},
	\end{equation}
	which is the energy of the ground state of $H$ in the space of states
	satisfying $\langle a(J)|A|a(J)\rangle = A_c$.

	We now translate this to field theory. In field theory the only
	normalizable eigenstate of a Hamiltonian is the vacuum, so $|a(J)\rangle$
	is identified with the vacuum of the sourced theory. Adding a source
	$J(x)\varphi(x)$ to the Hamiltonian density, adiabatically switched on
	at $t=0$ and off at $t=T$ over spatial volume $V$, the adiabatic theorem
	gives
	\begin{equation}
		e^{iW(J)} = \langle 0^+|0^-\rangle
		= \exp\left[-iVT\,\varepsilon(J)\right],
	\end{equation}
	so $-W(J)$ is the ground state action of the sourced theory. The Legendre
	transform structure,
	\begin{equation}
		\Gamma(\bar{\varphi}) = W(J) - J\frac{\delta W}{\delta J},
		\qquad \frac{\delta W}{\delta J} = \bar{\varphi},
	\end{equation}
	is formally identical to the quantum-mechanical result above, identifying
	$-\Gamma(\bar\varphi)$ as the ground state action of the sourceless
	theory in the sector with $\langle a|\varphi|a\rangle = \bar\varphi$.
	That is,
	\begin{align}
		-\Gamma(\bar{\varphi})            & = \langle a|H|a\rangle,            \\
		-\frac{\Gamma(\bar{\varphi})}{VT} & = \frac{\langle a|H|a\rangle}{VT},
	\end{align}
	and therefore
	\begin{equation}
		V_{\text{eff}}(\bar{\varphi}) = \langle a|\mathcal{H}|a\rangle
	\end{equation}
	for the lowest-energy state satisfying $\langle a|a\rangle=1$ and
	$\langle a|\varphi|a\rangle = \bar\varphi$.
\end{proof}

This result has an important consequence: since $V_{\text{eff}}(\bar\varphi)$
is defined as a minimum of energy over a convex set of states, it is
convex. Its minima therefore correspond to the stable equilibrium
configurations of the quantum field theory, a fact that will be central
to the analysis of phase transitions in later chapters.

\section{Variational Definition of the Full Effective Action}

The variational interpretation of $V_{\text{eff}}$ established above can
be generalized to the full effective action.

\begin{theorem}
	$\Gamma(\varphi)$ is the stationary value of
	\begin{equation}
		\Gamma(\varphi) = \int_{-\infty}^{\infty} dt\,
		\langle\psi_-,t|\,i\partial_t - H\,|\psi_+,t\rangle
	\end{equation}
	subject to the constraints
	\begin{equation}
		\langle\psi_-,t|\Phi(\mathbf{x})|\psi_+,t\rangle = \varphi(t,\mathbf{x}),
		\qquad
		\langle\psi_-,t|\psi_+,t\rangle = 1,
		\qquad
		\lim_{t\to\pm\infty}|\psi_\pm,t\rangle = |0\rangle.
	\end{equation}
\end{theorem}

\begin{proof}
	The proof follows the same structure as Theorem~2. Introducing Lagrange
	multipliers $J(t,\mathbf{x})$ and $w(t)$ for the two constraints,
	\begin{multline}
		\Gamma_m(\varphi) = \int dt\,\langle\psi_-,t|i\partial_t -
		H|\psi_+,t\rangle
		+ \int d^4x\,J(t,\mathbf{x})\left(\langle\psi_-,t|\Phi|\psi_+,t\rangle
		- \varphi(t,\mathbf{x})\right) \\
		- w(t)\left(\langle\psi_-,t|\psi_+,t\rangle - 1\right).
	\end{multline}
	Varying with respect to $\langle\psi_-,t|$, the last term drops out and
	since $\langle\delta\psi_-,t|$ is arbitrary,
	\begin{equation}
		\boxed{\left(i\partial_t - H
			+ \int d^3x\,J(t,\mathbf{x})\,\Phi(\mathbf{x})\right)|\psi_+,t\rangle
			= w(t)|\psi_+,t\rangle.}
	\end{equation}
	Varying with respect to $|\psi_+,t\rangle$ and integrating by parts,
	\begin{equation}
		\int dt\,\langle\psi_-,t|i\partial_t|\delta\psi_+,t\rangle
		= \left[\langle\psi_-,t|\delta\psi_+,t\rangle\right]_{-\infty}^{\infty}
		- \int dt\,i\left(\partial_t\langle\psi_-,t|\right)|\delta\psi_+,t\rangle,
	\end{equation}
	where the boundary term vanishes because
	$\lim_{t\to\pm\infty}|\psi_\pm,t\rangle = |0\rangle$. Hence
	\begin{equation}
		\int dt\,\langle\psi_-,t|i\partial_t|\delta\psi_+,t\rangle
		= -\int dt\,i\left(\partial_t\langle\psi_-,t|\right)|\delta\psi_+,t\rangle.
	\end{equation}
	Since $|\delta\psi_+,t\rangle$ is arbitrary,
	\begin{equation}
		\delta_{|\psi_+,t\rangle}\Gamma_m(\varphi) = \int dt\,\langle\psi_-,t|
		\left(-i\overleftarrow{\partial_t} + H
		+ \int d^3x\,J(\mathbf{x})\,\Phi(\mathbf{x})
		+ w(t)\right)|\delta\psi_+,t\rangle = 0,
	\end{equation}
	and taking the Hermitian conjugate,
	\begin{equation}
		\boxed{\left[i\partial_t - H
				+ \int d^3x\,J(\mathbf{x})\,\Phi(\mathbf{x})\right]|\psi_-,t\rangle
			= w^*(t)|\psi_-,t\rangle.}
	\end{equation}
	Defining the rescaled states
	\begin{align}
		|+,t\rangle & = e^{\,i\int_{-\infty}^{t}dt'\,w(t')}\,|\psi_+,t\rangle, \\
		|-,t\rangle & = e^{-i\int_{t}^{\infty}dt'\,w^*(t')}\,|\psi_-,t\rangle,
	\end{align}
	these are asymptotic ground states satisfying the sourced Schr\"odinger
	equation with no inhomogeneous term. Their overlap gives
	\begin{equation}
		\langle-,t|+,t\rangle = e^{iW(J)}
		= \exp\!\left[i\int_{-\infty}^{\infty}dt\,w(t)\right],
	\end{equation}
	from which
	\begin{equation}
		W(J) = \int_{-\infty}^{\infty}dt\,
		\langle\psi_-,t|\,i\partial_t - H
		+ \int d^3x\,J\cdot\Phi\,|\psi_+,t\rangle.
	\end{equation}
	As shown in Appendix~\ref{app:delta-W-delta-J},
	$\delta W(J)/\delta J = \varphi$, so the Legendre transform gives
	\begin{equation}
		W(J) - \int d^4x\,J\varphi = \Gamma(\varphi)
		= \int_{-\infty}^{\infty}dt\,
		\langle\psi_-,t|\,i\partial_t - H\,|\psi_+,t\rangle,
	\end{equation}
	which shows that $\Gamma(\varphi)$ is indeed the stationary value.
\end{proof}

The equation \eqref{eq:gamma-eom} has a further physical interpretation.
When the source vanishes, $\delta\Gamma/\delta\bar\varphi = 0$ becomes
the exact quantum equation of motion for $\varphi(\mathbf{x},t)$.
Expanding $\Gamma$ in powers of momentum about vanishing external momenta,
\begin{equation}
	\Gamma(\varphi) = \int d^4x\left[
		-V_{\text{eff}}(\varphi)
		+ \frac{1}{2}\left(\partial_\mu\varphi\,\partial^\mu\varphi\right)
		Z(\varphi) + \cdots\right],
\end{equation}
and renormalizing the field so that $Z(\varphi) = 1$, the equation of
motion becomes
\begin{equation}
	\square\,\varphi = -V'_{\text{eff}}(\varphi).
\end{equation}
In the static limit this is precisely the classical Klein-Gordon equation,
but with the classical potential replaced by the effective potential. Thus
quantum corrections to the dynamics are fully captured by $V_{\text{eff}}$,
which is yet another reason it is the central object of study.
